% MT046 / explanation — canonical editable content.
一条样本\(x=(x_1,\ldots,x_d)\)含\(d\)个测量值，\(y\)是类别0或1。Gaussian NB要回答：这组测量更像哪一类？本题三问依次是“作什么假设、用什么分布、怎样作决定”。


\textbf{训练时分别学哪些量。} 类别先验\(\pi_c=P(y=c)\)表示\(c\)类的比例；\(\mu_{cj}\)表示该类第\(j\)维的平均值；\(\sigma^2_{cj}\)表示这维围绕平均值的波动。方差大，分布较宽；方差小，同样偏离均值的距离会显得更异常。不同类通常各学一组均值、方差。


\textbf{两个不同层面的独立假设。} 独立同分布（i.i.d.）针对训练记录：每条记录按同一总体规律抽取，相互独立。朴素假设针对同一记录内的特征：给定类别后各维独立。因此可以把各维高斯密度相乘：



\[
p(x\mid c)=\prod_{j=1}^d\frac{1}{\sqrt{2\pi\sigma_{cj}^2}}
\exp\left[-\frac{(x_j-\mu_{cj})^2}{2\sigma_{cj}^2}\right].
\]


这里\(\prod\)表示逐项相乘，\(\exp(t)=e^t\)。离均值越远，平方距离越大，指数越负，密度越小。前面的系数使每个高斯密度下的总面积为1；密度是分布曲线在这个位置的高度。一个很窄区间内的概率，约等于这里的密度乘以区间宽度；整条曲线下的面积为1。因此单点密度本身不是“恰好等于这个数的概率”。


\textbf{预测时把证据和先验合起来。} 比较\(\pi_c p(x\mid c)\)即可，因为后验共同的分母\(p(x)\)不改变排序。实际常取自然对数。这里\(e\approx2.718\)，\(\log q=s\)表示\(q=e^s\)。log严格递增，不改变大小关系，并把一串乘积改成求和。先看单个特征：令\(a=2\pi\sigma_{cj}^2>0\)、\(t=-(x_j-\mu_{cj})^2/(2\sigma_{cj}^2)\)，它的密度就是\(a^{-1/2}e^t\)。根号分母写成负二分之一次方后，取log得到



\[
\log(a^{-1/2}e^t)
=\log(a^{-1/2})+\log(e^t)
=-\frac12\log a+t.
\]


这里分别用了“乘积的log等于各项log之和”“幂次可以移到log前面”及\(\log(e^t)=t\)。把\(a,t\)换回原式，再将各维相加、加上\(\log\pi_c\)，得到



\[
g_c(x)=\log\pi_c-\frac12\sum_{j=1}^d\log(2\pi\sigma_{cj}^2)
-\frac12\sum_{j=1}^d\frac{(x_j-\mu_{cj})^2}{\sigma_{cj}^2}.
\]


先验、分布宽窄、离中心的距离都在这个分数里。\(\arg\max_c g_c(x)\)的意思是“返回使分数最大的类别编号”，不是返回分数本身。错分代价相同才直接采用这个规则。


\textbf{把补充例子从头算完。} 设两类先验各0.5，均值\(\mu_0=(0,0)\)、\(\mu_1=(1,1)\)，两维方差均为1，新点\(x=(0.8,0.8)\)。两类的\(\log0.5\)相同，\(-\log(2\pi)\)也相同。剩下的平方距离为



\[
D_0=0.8^2+0.8^2=1.28,\qquad D_1=(-0.2)^2+(-0.2)^2=0.08.
\]


所以\(g_0=K-0.64\)、\(g_1=K-0.04\)，其中\(K=\log0.5-\log(2\pi)\)是共同项。类1高0.60，应选1。分数是未归一化概率的log，因此分数差0.60表示两类概率之比为\(e^{0.60}\approx1.822\)。把类0的相对份额记为1，类1就是1.822；两份加起来再归一化，得到\(P(1\mid x)=1.822/(1+1.822)\approx0.646\)。这里能只比较普通距离，是因为先验与方差都相同；一般情形要保留完整分数。


选读：矩阵写法把各维均值装成向量\(\mu_c\)，把方差放到协方差矩阵对角线上。非对角位置为0表示模型没有加入各维间的协方差。这个紧凑写法与上面的逐维乘积表示同一模型。


\textbf{出处：} 2023A期中 Q7，原卷第4页（10分）；2025A期中 Q11，原卷第8页（10分）


\textbf{答案依据：} 2023A期中Q7随卷解答；2025A Q11同题归并。
